% !TeX root = 01_tp2.tex
\documentclass[../main]{subfiles}
\begin{document}

\chapter{Normativas y Diseño asistido}
\img{images/manoalzada.png}{Dibujo Manual y croquis digital}
\section{Análisis y Bocetado Previo}
\actividad{Leer detenidamente cada consigna. Las respuestas deben
  estar justificadas y los dibujos realizados a mano alzada utilizando
regla, lápiz y respetando las proporciones.}

\begin{enumerate}
    % Preguntas de razonamiento y dibujo manual
  \item El diseño asistido por computadora requiere conocer
    previamente las normas de dibujo. ¿Por qué es crítico
    estandarizar (mediante normas IRAM, ISO o IEC) los símbolos
    electrónicos antes de digitalizarlos en un software?
  \item Investigue la normativa vigente y dibuje a mano alzada los
    símbolos esquemáticos de los siguientes componentes: Resistencia,
    Capacitor polarizado, Diodo LED y Transistor NPN..

    % Conectando el TP1 con el razonamiento físico
  \item En el trabajo anterior se estableció la diferencia entre el
    CAD de propósito general y el diseño ECAD.. Describa, con
    sus propias palabras y basándose en lo investigado, cuál es el
    flujo de trabajo lógico desde que se boceta un diagrama
    esquemático hasta que se fabrica la placa PCB (Placa de Circuito
    Impreso)

    % Análisis espacial sin necesidad de PC
  \item Para integrar un software de modelado 3D en un proyecto
    electrónico se pueden diseñar carcasas, paneles o soportes. Si
    tuviera que modelar la carcasa protectora para su proyecto en
    FreeCAD, ¿Qué medidas físicas de los componentes electrónicos
    reales necesitaría tomar previamente con un calibre en el taller?
    Enumere al menos cinco ejemplos de cotas críticas.

    % Ensayo para evitar IA
  \item \textbf{Escalas de Dibujo:} Si un microcontrolador real mide
    35 mm de largo por 15 mm de ancho, y necesita dibujarlo en su
    carpeta utilizando una escala de ampliación 2:1 para mostrar los
    detalles de sus pines, ¿Cuáles serán las medidas finales en su
    dibujo? Explique por qué es vital indicar la escala en el rótulo del plano.

  \item \textbf{Tipos de Líneas en Mecánica Electrónica:} Al dibujar
    el plano de la carcasa de una fuente de alimentación, ¿Qué
    diferencia conceptual existe entre utilizar una línea continua
    gruesa y una línea de trazos (punteada)? Dibuje un ejemplo simple
    de un cubo hueco aplicando ambas líneas.

  \item \textbf{Lectura de Planos (Sistemas):} Describa con sus
    propias palabras qué es un ``Diagrama de Bloques'' y en qué se
    diferencia de un ``Diagrama Esquemático'' detallado. Dibuje un
    diagrama de bloques secuencial de un sistema de audio básico
    (Micrófono, Preamplificador, Amplificador de Potencia y Altavoz).

  \item \textbf{Análisis de Fallas en Pistas (PCB):} Si al dibujar
    manualmente las pistas de un circuito impreso (routing) deja muy
    poco espacio libre (por ejemplo, menos de 0.5 mm) entre dos
    pistas de cobre, ¿Qué problema físico o eléctrico podría ocurrir
    al momento de realizar la soldadura manual con estaño?

  \item \textbf{Simbología y Polaridad:} Elija un componente
    electrónico real que tenga polaridad (por ejemplo, un capacitor
    electrolítico o un diodo). Dibuje a mano alzada su símbolo
    técnico normalizado e indique, mediante flechas, cómo se
    identifica físicamente esa misma polaridad en el componente real
    que usamos en el taller.

  \item Redacte una breve conclusión (mínimo 150 palabras) explicando
    cómo los programas de diseño asisten en la fabricación,
    pero por qué no pueden reemplazar el criterio del técnico al
    momento de distribuir físicamente los componentes por disipación
    de calor o interferencias.
    % --- Sección de Razonamiento Espacial y Normativa ---

\end{enumerate}

\end{document}
